\quad Em várias situações do nosso dia a dia, nos deparamos com a necessidade de trabalhar com dados e informações ordenados de forma correta, como, por exemplo, ao procurar um contato na lista telefônica. Seria difícil localizar um nome caso os contatos não estivessem organizados em ordem alfabética. Dessa forma, fica fácil notar que as atividades que envolvem algum método de ordenação estão muito presentes na computação.

\quad Os algoritmos de ordenação são responsáveis por organizar um conjunto de elementos de acordo com determinado critério, como ordem crescente, decrescente ou alfabética. Essa organização facilita a busca, a visualização e o processamento dos dados, contribuindo para a eficiência de diversas aplicações.

\quad Existem diferentes algoritmos de ordenação, cada um apresentando características, estratégias e custos computacionais próprios. Entre eles, destacam-se os algoritmos \textit{Insertion Sort}, \textit{Bubble Sort}, \textit{Selection Sort}, \textit{Shell Sort}, \textit{Merge Sort} e \textit{Quick Sort}. Apesar de possuírem o mesmo objetivo, esses algoritmos utilizam diferentes métodos para organizar os elementos de um vetor.

\quad O \textit{Insertion Sort}, também conhecido como ordenação por inserção, constrói progressivamente uma parte ordenada do vetor, inserindo cada novo elemento em sua posição correta. O \textit{Bubble Sort} realiza comparações entre elementos adjacentes e troca suas posições quando estão fora de ordem. Já o \textit{Selection Sort} procura o menor elemento da parte ainda não ordenada e o posiciona na posição correspondente. O \textit{Shell Sort} utiliza intervalos, chamados de \textit{gaps}, para realizar ordenações parciais. O \textit{Merge Sort} utiliza a estratégia de divisão e conquista, dividindo o vetor em partes menores e intercalando os resultados até obter a ordenação completa. Por sua vez, o \textit{Quick Sort} também adota a estratégia de divisão e conquista, mas realiza o particionamento do vetor em torno de um elemento escolhido, chamado de pivô, de modo que os elementos menores fiquem de um lado e os maiores do outro, repetindo o processo recursivamente em cada partição.

\quad No \textit{Quick Sort}, a escolha do pivô influencia diretamente o desempenho do algoritmo. Por esse motivo, foram implementadas duas versões: na primeira, o pivô é o primeiro elemento do vetor (ou da partição); na segunda, o pivô é definido a partir da média aritmética dos elementos da partição. Essa diferença pode alterar o balanceamento das partições e, consequentemente, o tempo de execução, principalmente em entradas já ordenadas ou em ordem decrescente.

\quad A análise da complexidade desses algoritmos é importante para compreender como o tempo de execução varia de acordo com o tamanho e a organização inicial da entrada. Dependendo do algoritmo e do caso analisado, podem existir diferenças significativas entre o melhor caso, o caso médio e o pior caso. No \textit{Shell Sort}, além do tamanho da entrada, o comportamento também está relacionado à sequência de \textit{gaps} utilizada. Já no \textit{Quick Sort}, o desempenho depende da qualidade da escolha do pivô, uma vez que partições desbalanceadas levam ao pior caso, enquanto partições equilibradas favorecem o melhor caso e o caso médio.

\quad Neste trabalho, serão apresentados o funcionamento e a implementação dos algoritmos \textit{Insertion Sort}, \textit{Bubble Sort}, \textit{Selection Sort}, \textit{Shell Sort}, \textit{Merge Sort} e \textit{Quick Sort} (nas versões com pivô no primeiro elemento e com pivô pela média aritmética), juntamente com testes de mesa e análises de complexidade. Também serão realizados experimentos utilizando entradas crescentes, decrescentes e aleatórias, considerando diferentes instâncias.